\bookchapter{Capítulo 43: Engenharia de diálogo — técnicas de conversa para colaborar com a IA com eficiência}

\chapterintro{O diálogo é a forma central de interação na programação com IA. Mas atenção: 'diálogo' não é a mesma coisa que 'conversa fiada' — um diálogo eficiente de programação com IA é uma técnica de comunicação estruturada, que chamamos de 'engenharia de diálogo'. Ela estuda como obter o código de maior qualidade com o menor número possível de rodadas de conversa. Simplificando, este capítulo ensina como deixar as coisas claras ao conversar com a IA.}

\booksection{Estratégia de camadas de Prompt}

O diálogo com a IA pode ser dividido em três camadas, e cada uma tem uma função diferente. Ao entender essas três camadas, você combina Prompts eficientes como quem empilha blocos de montar.

\begin{center}
\begin{tabularx}{\linewidth}{|X|X|X|X|}
\hline
\textbf{Camada} & \textbf{Função} & \textbf{Exemplo} & \textbf{Frequência} \tabularnewline
\hline
Camada de sistema & Define o papel da IA e as restrições globais & 'Você é um especialista em React e o código deve ser em TypeScript' & Configurada uma vez no início do projeto \tabularnewline
Camada de tarefa & Descreve a tarefa específica a fazer agora & 'Implemente um componente de formulário com campos de nome de usuário e senha' & Enviada no início de cada tarefa \tabularnewline
Camada de feedback & Avalia e corrige o que a IA produziu & 'O campo de senha está sem alternar mostrar/ocultar; adicione, por favor' & Enviada conforme a qualidade do resultado \tabularnewline
\hline
\end{tabularx}
\end{center}

\begin{center}\small\color{gray}As três camadas do Prompt de diálogo\end{center}

O Prompt da camada de sistema geralmente fica no arquivo de memória do projeto (como o CLAUDE.md) e não precisa ser repetido a cada conversa. A camada de tarefa é a que você mais usa na conversa — ela deve conter três elementos: 'o que fazer', 'com qual tecnologia' e 'quais restrições'. A camada de feedback é a que mais exige habilidade — um bom feedback permite que a IA corrija com precisão; um feedback ruim só a faz ficar chutando de novo.

\begin{tcbinfo}{Dica: os três golpes do feedback}
Um modelo prático de feedback é 'apontar o problema + explicar o esperado + fornecer contexto'. Por exemplo: 'o tratamento de erros da linha 25 usa apenas console.log (aponta o problema), em produção deveria usar um handler de erros unificado e devolver mensagens amigáveis ao usuário (explica o esperado), e o projeto já tem uma função handleError que pode ser usada (fornece o contexto)'. Esse tipo de feedback é dez vezes mais eficaz do que 'melhora o tratamento de erros'.

\end{tcbinfo}

\booksection{Diálogo de depuração: como descrever erros para a IA localizar rápido}

A depuração é um dos cenários mais frequentes na programação com IA. Quando você encontra um bug, o jeito de descrever o problema à IA determina diretamente se ela vai conseguir localizar e corrigir rápido. A seguir, os quatro pontos-chave do diálogo de depuração:

\begin{enumerate}
  \item Mostre à IA a informação original, e não o seu relato. Mencione diretamente com @ o log de erros, a pilha de exceções e o código relevante, em vez de descrever com um 'parece que deu algum erro'. Para a IA, ver a informação original é mais preciso do que ver a sua interpretação.
  \item Descreva a diferença entre o 'comportamento esperado' e o 'comportamento real'. 'Eu esperava que, ao clicar no botão, fosse para a página inicial (esperado), mas, na verdade, nada acontece ao clicar (real)' — esse tipo de descrição é muito mais útil do que 'o botão está quebrado'.
  \item Forneça os passos para reproduzir. 'Primeiro digito o nome de usuário, depois clico no botão de entrar, e o erro aparece' — os passos de reprodução permitem que a IA simule o seu caminho de operação e localize o problema mais facilmente.
  \item Diga claramente o que você já tentou. 'Já verifiquei a requisição de rede: retorna 200 e os dados estão normais' — isso ajuda a IA a descartar caminhos já verificados e a não perder tempo repetindo verificações.
\end{enumerate}

Descrição de depuração ruim

\begin{lstlisting}
O botão não responde ao clique, dá uma olhada aí
\end{lstlisting}

Descrição de depuração boa

\begin{lstlisting}
Após clicar no botão Entrar, a página não reage. Comportamento esperado: após o clique, ir para a página /home Comportamento real: nada acontece ao clicar e não há erros no console Passos para reproduzir: 1. Na página de login, digitar o usuário admin e a senha 123456 2. Clicar no botão 'Entrar' 3. A página não reage Já verificado: - Requisição de rede normal, retorna 200 - Evento onClick já vinculado - Código do componente LoginButton revisado Arquivos relacionados: @src/components/LoginButton.tsx @src/pages/Login.tsx
\end{lstlisting}

\booksection{Estratégias de iteração: passos pequenos e rápidos vs. refatoração em grandes passos}

A iteração na programação com IA tem duas estratégias básicas, para cenários diferentes. Entender a diferença entre elas ajuda você a escolher corretamente em cada situação.

A estratégia de 'passos pequenos e rápidos' serve para o desenvolvimento de funcionalidades e a correção de bugs. A essência: a cada vez, peça à IA apenas uma pequena alteração e só passe para a próxima depois de validar. Vantagens: cada passo é verificável e reversível, com risco controlado; desvantagens: é mais lento e exige interações frequentes de conversa. Ideal para: desenvolvimento de novas funcionalidades, exploração quando você não tem certeza da melhor solução e iteração cautelosa quando não confia muito no resultado da IA.

A estratégia de 'refatoração em grandes passos' serve para modificações em larga escala quando o objetivo está claro. A essência: descreva o objetivo completo de uma vez e deixe a IA planejar e executar toda a mudança de forma autônoma. Vantagens: é rápido e a IA enxerga o quadro geral, tomando decisões de arquitetura melhores; desvantagens: o risco é alto e, se a IA entender errado, corrigir custa caro. Ideal para: upgrades de arquitetura bem definidos, migrações de código em grande escala e refatorações com cobertura de testes sólida.

\begin{center}
\begin{tabularx}{\linewidth}{|X|X|X|}
\hline
\textbf{Dimensão} & \textbf{Passos pequenos e rápidos} & \textbf{Refatoração em grandes passos} \tabularnewline
\hline
Volume por alteração & Um ponto de funcionalidade ou um bug & Um módulo ou uma camada de arquitetura \tabularnewline
Frequência de validação & Valida imediatamente após cada alteração & Valida tudo junto, após concluir todas as alterações \tabularnewline
Rodadas de conversa & Muitas (5 a 10 por funcionalidade) & Poucas (1 a 3 por módulo) \tabularnewline
Risco & Baixo; cada passo é reversível & Alto; corrigir desvios custa caro \tabularnewline
Cenários adequados & Novas funcionalidades, desenvolvimento exploratório & Refatorações e migrações com objetivo claro \tabularnewline
Pré-requisitos & Nenhum requisito especial & Precisa de cobertura de testes \tabularnewline
\hline
\end{tabularx}
\end{center}

\begin{center}\small\color{gray}Comparação entre as duas estratégias de iteração\end{center}

Na prática, as duas estratégias costumam ser combinadas. Por exemplo, ao refatorar um módulo, primeiro use a 'refatoração em grandes passos' para a IA apresentar o plano geral; depois de revisar o plano, execute aos poucos com 'passos pequenos e rápidos'. Assim você mantém a visão global e controla o risco de execução.

\booksection{Gerenciamento de contexto em conversas de múltiplas rodadas}

Conforme o número de rodadas aumenta, você precisa gerenciar ativamente o contexto da conversa para evitar o problema de 'a IA esqueceu as informações do início'. A seguir, algumas técnicas práticas de gerenciamento de contexto.

Primeiro, 'ancore' as informações-chave. No início da conversa, deixe as restrições importantes claras por escrito e reafirme-as nos momentos certos ao longo do diálogo. Por exemplo: 'lembre-se: nossos dados ficam no localStorage, não podemos usar banco de dados' — reitere essa restrição nas solicitações de alteração relacionadas, para evitar que a IA esqueça.

Segundo, use bem a técnica de 'resumir e continuar'. Quando a conversa ficar muito longa, peça à IA que 'resuma o progresso atual e as pendências' e então abra uma nova conversa tomando esse resumo como ponto de partida. Assim você preserva o contexto essencial e evita a degradação das conversas longas.

Terceiro, 'uma conversa, uma tarefa'. Essa é a estratégia de gerenciamento de contexto mais simples e mais eficaz. Cada conversa trata apenas de uma tarefa ou de um ponto de funcionalidade; ao concluir, abra uma nova conversa para o próximo. Isso limita naturalmente o comprimento da conversa e mantém o contexto claro.

Quarto, aproveite os recursos de gerenciamento de contexto das ferramentas. A maioria das ferramentas de programação com IA tem botões de 'limpar contexto' ou 'nova sessão' — quando sentir que a qualidade das respostas da IA piorou, não hesite: abra uma nova sessão. Perder parte do histórico da conversa é muito melhor do que continuar sofrendo com um contexto degradado.

\begin{tcbpractice}{Mãos à obra: seja o 'paciente' da sua IA uma vez}
Monte de propósito uma pequena página no navegador com um bug (por exemplo, um botão que não reage). Depois, descreva tudo à IA uma vez, seguindo o formato de 'boa descrição de depuração' deste capítulo: informação original, esperado vs. real, passos de reprodução e o que já foi verificado. Sinta a diferença na velocidade com que a IA localiza o problema — você vai entender por que 'saber descrever o problema' é, em si, uma habilidade.

\end{tcbpractice}

\begin{tcbsummary}{Resumo do capítulo}
\begin{itemize}
  \item Três camadas do Prompt: sistema (vai no arquivo de memória), tarefa (o que fazer + tecnologia + restrições) e feedback (problema + esperado + contexto)
  \item Quatro pontos da depuração: mostrar a informação original, esperado vs. real, passos de reprodução e o que já foi verificado
  \item Passos pequenos e rápidos: risco baixo e muitas rodadas, ideal para novas funcionalidades; refatoração em grandes passos: risco alto e poucas rodadas, exige cobertura de testes
  \item Quatro truques de contexto: ancorar as informações-chave, resumir e continuar, uma conversa por tarefa e abrir uma nova sessão sem hesitar
\end{itemize}

\end{tcbsummary}
